\chapter{Logic, Comparisons, and Strings}
\label{chap:logic-strings}

\chapterguide
  {You should be able to create variables and print results.}
  {use logical and relational operators; create valid strings; and concatenate, format, measure, transform, and extract text.}

\section{Logical operations}

The core logical operators are:

\begin{dstable}
\begin{tabularx}{\linewidth}{@{}p{2.2cm}X@{}}
\toprule
\rowcolor{DSPaleBlue!92}
Operator & Meaning in the handout \\
\midrule
\texttt{!} & Logical NOT \\
\texttt{\&} & Logical AND \\
\texttt{|} & Logical OR \\
\bottomrule
\end{tabularx}
\end{dstable}

With single Boolean values:

\begin{lstlisting}[style=dsr]
a = TRUE
b = FALSE

print(a & b)
print(a & !b)
print(a | b)
print(!a | b)
\end{lstlisting}

The following code compares element-wise and short-circuit forms:

\begin{lstlisting}[style=dsr]
c = c(TRUE, FALSE)
d = c(FALSE, FALSE)

print(c && d)
print(c || d)

v <- c(3, 1.5, TRUE)
t <- c(4, 2.1, FALSE)

print(!v)
print(!t)
print(v & t)
print(v | t)
print(v && t)
print(v || t)
\end{lstlisting}

\begin{keyidea}
\texttt{\&} and \texttt{|} operate element by element. \texttt{\&\&} and \texttt{||} inspect only the first element and may skip evaluation of the second expression. Use the double-character forms for scalar control-flow conditions.
\end{keyidea}

\begin{utpfig}[Element-wise versus short-circuit]{The same two vectors, read by
two different operators. \texttt{\&} pairs every position and returns a vector
as long as its inputs. \texttt{\&\&} reads position~1 only and returns a single
value, which is why the double forms belong in \texttt{if} conditions --- those
expect exactly one \texttt{TRUE} or \texttt{FALSE}.}{fig:logic-elementwise}
\begin{tikzpicture}[x=1cm,y=1cm]

% ---- left panel: element-wise & ---------------------------------
\node[tagb] at (1.35,1.15) {\texttt{c \& d}};
\node[tag]  at (1.35,0.78) {every position};

\node[tag,anchor=east] at (-0.10,0) {\texttt{c}};
\node[cell] at (0.68,0) {TRUE};
\node[cell] at (2.02,0) {FALSE};

\node[tagb,anchor=east] at (-0.10,-0.85) {\texttt{\&}};
\node[cell] at (0.68,-0.85) {FALSE};
\node[cell] at (2.02,-0.85) {FALSE};

\draw[semithick,draw=UTPMuted!65] (0.16,-1.30) -- (2.54,-1.30);

\node[cellout] at (0.68,-1.80) {FALSE};
\node[cellout] at (2.02,-1.80) {FALSE};
\node[tag] at (1.35,-2.38) {a vector of length 2};

% ---- right panel: short-circuit && ------------------------------
\node[tagb] at (5.75,1.15) {\texttt{c \&\& d}};
\node[tag]  at (5.75,0.78) {position 1 only};

\node[tag,anchor=east] at (4.30,0) {\texttt{c}};
\node[cell]      at (5.08,0) {TRUE};
\node[cellghost] at (6.42,0) {FALSE};

\node[tagb,anchor=east] at (4.30,-0.85) {\texttt{\&\&}};
\node[cell]      at (5.08,-0.85) {FALSE};
\node[cellghost] at (6.42,-0.85) {FALSE};

\draw[semithick,draw=UTPMuted!65] (4.56,-1.30) -- (5.60,-1.30);

\node[cellout] at (5.08,-1.80) {FALSE};
\node[tag] at (5.08,-2.38) {one value};
\node[tag,text width=1.7cm,align=center] at (6.48,-1.80) {not evaluated};

\end{tikzpicture}
\end{utpfig}

\section{Relational operations}

Relational operators compare values and return logical results.

\begin{dstable}
\begin{tabularx}{\linewidth}{@{}p{2.4cm}X@{}}
\toprule
\rowcolor{DSPaleBlue!92}
Operator & Meaning \\
\midrule
\texttt{<} & Less than \\
\texttt{>} & Greater than \\
\texttt{<=} & Less than or equal to \\
\texttt{>=} & Greater than or equal to \\
\texttt{==} & Equal to \\
\texttt{!=} & Not equal to \\
\bottomrule
\end{tabularx}
\end{dstable}

\begin{lstlisting}[style=dsr]
x = 5
y = 9
z = 2 + 3

print(x > y)
print(x < y)
print(x <= z)
print(y >= z)
print(x == z)
print(x != y)
\end{lstlisting}

Because \texttt{z} evaluates to \texttt{5}, \texttt{x == z} is TRUE.

Comparisons also operate element by element on vectors:

\begin{lstlisting}[style=dsr]
v <- c(2, 5.5, 6, 9)
t <- c(8, 2.5, 14, 9)

print(v > t)
print(v < t)
print(v <= t)
print(v >= t)
print(v == t)
print(v != t)
\end{lstlisting}

\section{Constructing strings}

R strings may use single or double quotation marks:

\begin{lstlisting}[style=dsr]
a <- 'Start and end with single quote'
b <- "Start and end with double quotes"
c <- "single quote ' in between double quotes"
d <- 'Double quotes " in between single quote'
\end{lstlisting}

Embedded delimiters of the same type must be escaped.

\begin{pitfall}
Unmatched or unescaped delimiters cause syntax errors. Keep quotation marks balanced.
\end{pitfall}

\section{Escape characters}

Common escape sequences include:

\begin{dstable}
\begin{tabularx}{\linewidth}{@{}p{2.4cm}X@{}}
\toprule
\rowcolor{DSPaleBlue!92}
Sequence & Description \\
\midrule
\texttt{\textbackslash"} & Double quote \\
\texttt{\textbackslash\textbackslash} & Backslash \\
\texttt{\textbackslash n} & New line \\
\texttt{\textbackslash r} & Carriage return \\
\texttt{\textbackslash t} & Tab \\
\bottomrule
\end{tabularx}
\end{dstable}

\begin{lstlisting}[style=dsr]
str <- "We are the so-called \"Vikings\", from the north."
cat(str)

str <- "We are the so-called \\Vikings\\, from the north."
cat(str)

str <- "We are the \rso-called Vikings\nfrom the north."
cat(str)

str <- "\tWe are the so-called Vikings from the north."
cat(str)
\end{lstlisting}

\texttt{cat()} renders the escape sequences instead of printing a quoted representation.

\section{Concatenating strings}

\texttt{paste()} combines text with configurable separators:

\begin{lstlisting}[style=dsr]
a <- "Hello"
b <- "How"
c <- "are you? "

print(paste(a, b, c))
print(paste(a, b, c, sep = "-"))
print(paste(a, b, c, sep = "", collapse = ""))
\end{lstlisting}

\texttt{sep} joins arguments within each result; \texttt{collapse} joins multiple result elements.

\begin{utpfig}[\texttt{sep} joins across, \texttt{collapse} joins down]{The two
arguments answer different questions. \texttt{sep} sits between the pieces that
form one result element; \texttt{collapse} sits between the result elements
themselves. Supply \texttt{sep} alone and the result stays a vector. Supply
\texttt{collapse} as well and it becomes a single string.}{fig:paste-sep-collapse}
\begin{tikzpicture}[x=1cm,y=1cm]

\node[tag,anchor=east] at (-0.10,0)     {\texttt{x}};
\node[tag,anchor=east] at (-0.10,-0.80) {\texttt{y}};

\node[cell]    at (0.70,0)     {\texttt{"A"}};
\node[cell]    at (0.70,-0.80) {\texttt{"1"}};
\node[cellalt] at (2.05,0)     {\texttt{"B"}};
\node[cellalt] at (2.05,-0.80) {\texttt{"2"}};

\draw[decorate,decoration={brace,amplitude=3pt},draw=UTPBlue,semithick]
  (0.28,0.36) -- (0.28,-1.16);
\draw[decorate,decoration={brace,amplitude=3pt},draw=UTPOrange,semithick]
  (1.63,0.36) -- (1.63,-1.16);

\node[tag,text width=2.4cm,align=center] at (1.38,-1.62)
  {\texttt{sep} joins within a position};

% first result: a vector of two strings
\node[tagb,anchor=west] at (4.05,0.62) {\texttt{paste(x, y, sep = "-")}};
\node[cellout] at (4.95,0)     {\texttt{"A-1"}};
\node[cellout] at (6.35,0)     {\texttt{"B-2"}};
\node[tag,anchor=west] at (7.10,0) {a vector of 2};
\draw[rarrow] (2.80,0) -- (4.20,0);

% second result: collapsed to one string
\draw[rarrow] (5.65,-0.42) -- (5.65,-1.06);
\node[tagb] at (5.65,-1.34) {\texttt{..., collapse = "+"}};
\node[cellout,minimum width=22mm] at (5.65,-1.94) {\texttt{"A-1+B-2"}};
\node[tag,anchor=west] at (7.05,-1.94) {one string};

\end{tikzpicture}
\end{utpfig}

\section{Formatting numbers and strings}

\texttt{format()} controls display without changing the underlying numeric value:

\begin{lstlisting}[style=dsr]
# Total number of displayed digits; final digit is rounded.
result <- format(23.123456789, digits = 9)
print(result)

# Scientific notation.
result <- format(c(6, 13.14521), scientific = TRUE)
print(result)

# Minimum digits to the right of decimal point.
result <- format(23.47, nsmall = 5)
print(result)

# format() returns formatted text.
result <- format(6)
print(result)

# Width and justification.
result <- format(13.7, width = 6)
print(result)
result <- format("Hello", width = 8, justify = "l")
print(result)
result <- format("Hello", width = 8, justify = "c")
print(result)
\end{lstlisting}

\section{Measuring and changing strings}

Use \texttt{nchar()} in base R or \texttt{str\_length()} from \texttt{stringr} to count characters:

\begin{lstlisting}[style=dsr]
str <- "Count the number of characters"
print(nchar(str))

install.packages("stringr")
library(stringr)
str_length(str)
\end{lstlisting}

Case conversion:

\begin{lstlisting}[style=dsr]
print(toupper("Changing To Upper"))
print(tolower("Changing To Lower"))
\end{lstlisting}

Case conversion supports case-insensitive comparison and standardized output.

\section{Extracting parts of a string}

\texttt{substring()} and \texttt{substr()} extract character ranges:

\begin{lstlisting}[style=dsr]
# 5th through 7th characters
result <- substring("Extract", 5, 7)
print(result)

# first character
result <- substr("Learn Code Tech", 1, 1)
print(result)
\end{lstlisting}

\begin{utpfig}[Character positions in a string]{\texttt{substring()} and
\texttt{substr()} count characters from~1, and both endpoints are included. The
call keeps positions 5, 6 and 7, so seven characters become three. Counting
from~0, as several other languages do, is the usual source of an off-by-one
result here.}{fig:substring}
\begin{tikzpicture}[x=1cm,y=1cm]

\foreach \ch/\i in {E/1, x/2, t/3, r/4} {
  \node[cell] at (\i*0.92,0) {\texttt{\ch}};
  \node[tag]  at (\i*0.92,-0.62) {\i};
}
\foreach \ch/\i in {a/5, c/6, t/7} {
  \node[cellout] at (\i*0.92,0) {\texttt{\ch}};
  \node[tagb]    at (\i*0.92,-0.62) {\i};
}

\node[tag,anchor=east] at (0.46,0) {\texttt{"Extract"}};

\draw[decorate,decoration={brace,amplitude=4pt,mirror},draw=UTPTeal,semithick]
  (4.28,-0.90) -- (6.76,-0.90);

\node[tagb,anchor=north] at (5.52,-1.16) {\texttt{substring("Extract", 5, 7)}};

\node[cellout] at (8.55,0) {\texttt{"act"}};
\draw[rarrow] (7.28,0) -- (8.05,0);

\end{tikzpicture}
\end{utpfig}

\begin{quickcheck}
\begin{enumerate}
  \item What type of value is returned by a relational comparison such as \texttt{x < y}?
  \item Why does \texttt{x == z} evaluate as true in the handout example?
  \item What is the role of a backslash in \texttt{\textbackslash"Vikings\textbackslash"} inside a double-quoted string?
  \item Which function in the handout converts text to uppercase?
  \item Which functions are used to extract character ranges from a string?
\end{enumerate}
\end{quickcheck}

\begin{chapterrecap}
\begin{itemize}
  \item Use element-wise logical operators for vectors and short-circuit forms for scalar conditions.
  \item Strings require balanced delimiters; escape sequences represent special characters.
  \item Base R provides functions for joining, formatting, measuring, changing, and extracting text.
\end{itemize}
\end{chapterrecap}
